\section{Normal spaces}

\import{topology/separation.tex}
\import{set/choice.tex}

Normality here includes the topological-space domain but not \teeone.
The phrase ``normal space'' is a natural-language alias, while \teefour adds
\teeone explicitly.

\begin{definition}\label{is_normal}
    $X$ is normal iff $X$ is a topological space and
        for all $C,D\in\closeds{X}$ such that
            $C$ is disjoint from $D$ there exist $U,V\in\opens[X]$ such that
            $C\subseteq U$ and $D\subseteq V$ and $U$ is disjoint from $V$.
\end{definition}

\begin{classification}\label{normal_is_topological}
    Suppose $X$ is normal.
    Then $X$ is a topological space.
\end{classification}
\begin{proof}
    Follows by \obv{is_normal}.
\end{proof}

\begin{abbreviation}\label{normal_space}
    $X$ is a normal space iff $X$ is normal.
\end{abbreviation}

\begin{classification}\label{normal_space_is_topological}
    Let $X$ be a normal space.
    Then $X$ is a topological space.
\end{classification}
\begin{proof}
    Follows by \obv{is_normal}.
\end{proof}

\begin{abbreviation}\label{teefour}
    $X$ is \teefour\ iff $X$ is normal and $X$ is \teeone.
\end{abbreviation}

\begin{abbreviation}\label{teefour_space}
    $X$ is a \teefour-space iff $X$ is \teefour.
\end{abbreviation}

\begin{classification}\label{teefour_space_is_topological}
    Let $X$ be a \teefour-space.
    Then $X$ is a topological space.
\end{classification}
\begin{proof}
    Follows by \obv{is_normal}.
\end{proof}

\begin{proposition}[Shrinking an open neighbourhood]\label{normal_implies_closed_shrinking}
    Let $X$ be a normal space.
    Let $C\in\closeds{X}$.
    Let $U\in\opens[X]$.
    Suppose $C\subseteq U$.
    Then there exists $V\in\opens[X]$ such that
        $C\subseteq V\subseteq\closure{V}{X}\subseteq U$.
\end{proposition}
\begin{proof}
    $X$ is a topological space by notion.
    $U\subseteq\carrier[X]$ by \obv{opens_type}.
    Let $D=\carrier[X]\setminus U$.
    $D\in\closeds{X}$ by \auto{complement_of_open_elem_closeds}.
    $C$ is disjoint from $D$ by \obv{disjoint,setminus,subseteq}.
    Take $V,W\in\opens[X]$ such that
        $C\subseteq V$ and $D\subseteq W$ and
        $V$ is disjoint from $W$ by \obv{is_normal}.
    $V,W\subseteq\carrier[X]$ by \obv{opens_type}.
    Hence $V\subseteq\carrier[X]\setminus W$
        by \auto{subseteq,setminus,disjoint}.
    $\carrier[X]\setminus W\in\closeds{X}$
        by \obv{complement_of_open_elem_closeds}.
    \begin{align*}
        \closure{V}{X}
            &\subseteq\carrier[X]\setminus W
                \explanation{by \obv{closure_is_minimal_closed_set}}
            \\
            &\subseteq\carrier[X]\setminus D
                \explanation{by \obv{subseteq_implies_setminus_supseteq}}
            \\
            &=U
                \explanation{by \auto{double_relative_complement}}
    \end{align*}
    $V\subseteq\closure{V}{X}$
        by \obv{set_is_subseteq_to_closure_of_the_set}.
\end{proof}

\begin{proposition}\label{closed_shrinking_implies_normal}
    Let $X$ be a topological space.
    Suppose for all $C,U$ such that $C\in\closeds{X}$ and
        $U\in\opens[X]$ and $C\subseteq U$ there exists
        $V\in\opens[X]$ such that
        $C\subseteq V\subseteq\closure{V}{X}\subseteq U$.
    Then $X$ is normal.
\end{proposition}
\begin{proof}
    Show that for all $C,D\in\closeds{X}$ such that
        $C$ is disjoint from $D$ there exist $V,W\in\opens[X]$ such that
        $C\subseteq V$ and $D\subseteq W$ and
        $V$ is disjoint from $W$.
    \begin{subproof}
        Fix $C,D\in\closeds{X}$.
        Assume $C$ is disjoint from $D$.
        $C,D\subseteq\carrier[X]$ by \auto{closeds,pow_iff,subseteq}.
        Let $U=\carrier[X]\setminus D$.
        $U\in\opens[X]$ by \obv{closeds,is_closed_in}.
        Hence $C\subseteq U$
            by \auto{subseteq,setminus,disjoint}.
        Take $V\in\opens[X]$ such that
            $C\subseteq V\subseteq\closure{V}{X}\subseteq U$.
        Let $W=\carrier[X]\setminus\closure{V}{X}$.
        $\closure{V}{X}$ is closed in $X$
            by \obv{closure_is_closed}.
        $\closure{V}{X}\subseteq\carrier[X]$
            by \auto{closure_interior_frontier_is_in_carrier,opens_type}.
        Hence $W$ is open in $X$
            by \auto{complement_of_closed_is_open}.
        Thus $W\in\opens[X]$ by \obv{is_open}.
        $D\subseteq W$ by \auto{setminus,subseteq}.
        $V$ is disjoint from $W$
            by \obv{disjoint,setminus,subseteq}.
        Follows by \auto.
    \end{subproof}
    Follows by \obv{is_normal}.
\end{proof}

\begin{theorem}\label{normal_iff_closed_shrinking}
    Let $X$ be a topological space.
    Then $X$ is normal iff for all $C,U$ such that
        $C\in\closeds{X}$ and $U\in\opens[X]$ and $C\subseteq U$
        there exists $V\in\opens[X]$ such that
        $C\subseteq V\subseteq\closure{V}{X}\subseteq U$.
\end{theorem}
\begin{proof}
    Follows by
        \obv{normal_implies_closed_shrinking,closed_shrinking_implies_normal}.
\end{proof}

\begin{corollary}[Inserting an open set]\label{normal_insertion}
    Let $X$ be a normal space.
    Let $N\in\opens[X]$.
    Suppose $\closure{M}{X}\subseteq N$.
    Then there exists $L\in\opens[X]$ such that
        $\closure{M}{X}\subseteq L\subseteq\closure{L}{X}\subseteq N$.
\end{corollary}
\begin{proof}
    $\closure{M}{X}\in\closeds{X}$ by \obv{closure_elem_closeds}.
    Follows by \obv{normal_implies_closed_shrinking}.
\end{proof}

\begin{theorem}\label{teefour_space_is_teethree_space}
    Let $X$ be a \teefour-space.
    Then $X$ is a \teethree-space.
\end{theorem}
\begin{proof}
    $X$ is a topological space and $X$ is normal and $X$ is \teeone.
    Hence $X$ is a \teeone-space.
    Hence $X$ is a \teezero-space
        by \obv{teeone_space_is_teezero_space}.
    Thus $X$ is \teezero.
    Show that $X$ is regular.
    \begin{subproof}
        Show that for all $C,p$ such that
            $p\in\carrier[X]$ and $p\notin C\in\closeds{X}$
            there exist $U,V\in\opens[X]$ such that
            $p\in U$ and $C\subseteq V$ and $U\inter V=\emptyset$.
        \begin{subproof}
            Fix $C,p$.
            Assume $p\in\carrier[X]$ and $p\notin C$ and
                $C\in\closeds{X}$.
            $\{p\}$ is closed in $X$
                by \obv{teeone_implies_singletons_closed}.
            Hence $\{p\}\in\closeds{X}$ by \obv{is_closed_in}.
            $\{p\}$ is disjoint from $C$
                by \auto{disjoint,singleton_iff}.
            Take $U,V\in\opens[X]$ such that
                $\{p\}\subseteq U$ and $C\subseteq V$ and
                $U$ is disjoint from $V$ by \obv{is_normal}.
            Hence $p\in U$ by \auto{subseteq,singleton_iff}.
            Show that $U\inter V\subseteq\emptyset$.
            \begin{subproof}
                It suffices to show that every element of $U\inter V$
                    is an element of $\emptyset$ by \obv{subseteq}.
                Fix $x\in U\inter V$.
                Then $x\in U$ and $x\in V$ by \obv{inter}.
                Contradiction by \auto{disjoint}.
            \end{subproof}
            $\emptyset\subseteq U\inter V$ by \obv{emptyset_subseteq}.
            Hence $U\inter V=\emptyset$
                by \obv{subseteq_antisymmetric}.
            Follows by \auto.
        \end{subproof}
        Follows by \obv{is_regular}.
    \end{subproof}
    Hence $X$ is \teethree.
\end{proof}

\subsection{Normal shrinkers}

A normal shrinker chooses, for each closed set $C$ contained in an open set
$U$, an open neighbourhood of $C$ whose closure lies in $U$.
The pair law in \ref{normal_shrinker_apply_pair} expresses this property.
In a normal space, one application of Choice supplies all these neighbourhoods
at once; the Urysohn construction then iterates a fixed shrinker.

\begin{definition}\label{shrink_pairs}
    $\shrinkPairs{X}=\{p\in\closeds{X}\times\opens[X]\mid
        \fst{p}\subseteq\snd{p}\}$.
\end{definition}

\begin{proposition}\label{shrink_pair_elim}
    Let $p\in\shrinkPairs{X}$.
    Then $p=(\fst{p},\snd{p})$ and
        $\fst{p}\in\closeds{X}$ and $\snd{p}\in\opens[X]$ and
        $\fst{p}\subseteq\snd{p}$.
\end{proposition}
\begin{proof}
    $p\in\closeds{X}\times\opens[X]$ and
        $\fst{p}\subseteq\snd{p}$ by \obv{shrink_pairs}.
    Follows by \auto{times_elem_is_tuple,fst_type,snd_type}.
\end{proof}

\begin{proposition}\label{shrink_pair_intro}
    Let $C\in\closeds{X}$.
    Let $U\in\opens[X]$.
    Suppose $C\subseteq U$.
    Then $(C,U)\in\shrinkPairs{X}$.
\end{proposition}
\begin{proof}
    $(C,U)\in\closeds{X}\times\opens[X]$
        by \obv{times_tuple_intro}.
    $\fst{(C,U)}=C$ and $\snd{(C,U)}=U$ by \obv{fst_eq,snd_eq}.
    Follows by \obv{shrink_pairs}.
\end{proof}

\begin{definition}\label{shrink_relation}
    Let $X$ be a topological space.
    $\shrinkRelation{X}=\{q\in\shrinkPairs{X}\times\opens[X]\mid
        \fst{\fst{q}}\subseteq\snd{q}\subseteq
        \closure{\snd{q}}{X}\subseteq\snd{\fst{q}}\}$.
\end{definition}

\begin{proposition}\label{shrink_relation_is_relation}
    Let $X$ be a topological space.
    $\shrinkRelation{X}$ is a relation.
\end{proposition}
\begin{proof}
    $\shrinkRelation{X}\subseteq\shrinkPairs{X}\times\opens[X]$
        by \obv{shrink_relation,subseteq}.
    Follows by \auto{relation,subseteq,times_elem_is_tuple}.
\end{proof}

\begin{definition}\label{normal_shrinker}
    Let $X$ be a topological space.
    $h$ is a normal shrinker for $X$ iff
        $h\in\funs{\shrinkPairs{X}}{\opens[X]}$ and
        for every $p\in\shrinkPairs{X}$ we have
        $\fst{p}\subseteq h(p)\subseteq\closure{h(p)}{X}
            \subseteq\snd{p}$.
\end{definition}

\begin{proposition}\label{normal_shrinker_type}
    Let $X$ be a topological space.
    Let $h$ be a normal shrinker for $X$.
    Then $h\in\funs{\shrinkPairs{X}}{\opens[X]}$.
\end{proposition}
\begin{proof}
    Follows by \obv{normal_shrinker}.
\end{proof}

\begin{proposition}\label{normal_shrinker_apply}
    Let $X$ be a topological space.
    Let $h$ be a normal shrinker for $X$.
    Let $p\in\shrinkPairs{X}$.
    Then $h(p)\in\opens[X]$ and
        $\fst{p}\subseteq h(p)\subseteq\closure{h(p)}{X}
            \subseteq\snd{p}$.
\end{proposition}
\begin{proof}
    $h(p)\in\opens[X]$ by
        \auto{normal_shrinker_type,funs_type_apply}.
    Follows by \obv{normal_shrinker}.
\end{proof}

\begin{proposition}[The shrinking law]\label{normal_shrinker_apply_pair}
    Let $X$ be a topological space.
    Let $h$ be a normal shrinker for $X$.
    Let $C\in\closeds{X}$.
    Let $U\in\opens[X]$.
    Suppose $C\subseteq U$.
    Then $h((C,U))\in\opens[X]$ and
        $C\subseteq h((C,U))\subseteq\closure{h((C,U))}{X}
            \subseteq U$.
\end{proposition}
\begin{proof}
    $(C,U)\in\shrinkPairs{X}$ by \obv{shrink_pair_intro}.
    $\fst{(C,U)}=C$ and $\snd{(C,U)}=U$ by \obv{fst_eq,snd_eq}.
    Follows by \auto{normal_shrinker_apply}.
\end{proof}

\begin{proposition}\label{normal_shrink_relation_is_total}
    Let $X$ be a normal space.
    Then for every $p\in\shrinkPairs{X}$ there exists $V$ such that
        $p\mathrel{\shrinkRelation{X}}V$.
\end{proposition}
\begin{proof}
    Fix $p\in\shrinkPairs{X}$.
    $\fst{p}\in\closeds{X}$ and $\snd{p}\in\opens[X]$ and
        $\fst{p}\subseteq\snd{p}$ by \auto{shrink_pair_elim}.
    Take $V\in\opens[X]$ such that
        $\fst{p}\subseteq V\subseteq\closure{V}{X}\subseteq\snd{p}$
        by \obv{normal_implies_closed_shrinking}.
    $(p,V)\in\shrinkPairs{X}\times\opens[X]$
        by \obv{times_tuple_intro}.
    $\fst{(p,V)}=p$ and $\snd{(p,V)}=V$ by \obv{fst_eq,snd_eq}.
    Hence $(p,V)\in\shrinkRelation{X}$ by \obv{shrink_relation}.
\end{proof}

\begin{theorem}\label{normal_shrinker_exists}
    Let $X$ be a normal space.
    Then there exists $h$ such that $h$ is a normal shrinker for $X$.
\end{theorem}
\begin{proof}
    Let $S=\shrinkPairs{X}$.
    Let $R=\shrinkRelation{X}$.
    $R$ is a relation by \obv{shrink_relation_is_relation}.
    For every $p\in S$ there exists $V$ such that $p\mathrel{R}V$
        by \auto{normal_shrink_relation_is_total}.
    Take $h$ such that $h$ is a function on $S$ and
        for every $p\in S$ we have $p\mathrel{R}h(p)$
        by \obv{axiom_of_choice}.
    $h$ is a function and $S=\dom{h}$.
    Show that $h$ is a function to $\opens[X]$.
    \begin{subproof}
        Show that for every $p\in\dom{h}$ we have $h(p)\in\opens[X]$.
        \begin{subproof}
            Fix $p\in\dom{h}$.
            Hence $p\in S$.
            $(p,h(p))\in R$.
            Thus $(p,h(p))\in S\times\opens[X]$
                by \auto{shrink_relation}.
            Follows by \obv{times_tuple_elim}.
        \end{subproof}
    \end{subproof}
    Hence $h\in\funs{S}{\opens[X]}$ by \obv{funs_intro}.
    Show that for every $p\in S$ we have
        $\fst{p}\subseteq h(p)\subseteq\closure{h(p)}{X}
            \subseteq\snd{p}$.
    \begin{subproof}
        Fix $p\in S$.
        $(p,h(p))\in R$.
        $\fst{(p,h(p))}=p$ and $\snd{(p,h(p))}=h(p)$
            by \obv{fst_eq,snd_eq}.
        Follows by \auto{shrink_relation}.
    \end{subproof}
    $h$ is a normal shrinker for $X$ by \obv{normal_shrinker}.
\end{proof}
